
% !TeX spellcheck = pt_BR

\chapter{Sobre as ilustrações} \label{apendice-ilustrações}

Este livro foi dividido em três Atos, nos quais o autor enxergou que toda a trama se dividia. Cada um deles se abre com uma ilustração. Cada imagem foi concebida para
carregar, em linguagem visual, o mesmo argumento que o texto desenvolve
em linguagem matemática. São gravuras no estilo das ilustrações
científicas dos séculos XVII e XVIII --- época em que a Física ainda
precisava justificar a si mesma ao mundo.

O leitor que chega até este apêndice já percorreu os três Atos. Talvez,
ao reler agora as descrições abaixo, encontre nas imagens um significado
que na primeira vez ainda não estava disponível. É, afinal, o mesmo
mecanismo que governa toda boa compreensão: o que se vê depende
inteiramente do que já se sabe.


\section*{Primeiro Ato --- Cinemática}
% ─────────────────────────────────────────────────────────────────────



\begin{center}
    \textit{De Motu Projectilium} \\[0.2cm]
    {\small Sobre o movimento dos projéteis --- e sobre aprender a ver antes de aprender a explicar}
\end{center}

\bigskip

Há uma cena que gosto de imaginar. Ela se passa no século XVII, e um homem senta diante de um canhão disparado. Não
foge. Não se cobre. Ele \textit{observa}. Acompanha, com os olhos e com
a mente, o arco silencioso que a bala descreve no ar --- sobe, curva,
cai --- e se pergunta, talvez pela primeira vez na história, com rigor
matemático: \textit{qual é a forma exata dessa curva?}

O homem é Galileu Galilei, e a resposta que ele encontra muda o mundo.

A curva é uma parábola. Não aproximadamente --- exatamente. E Galileu a
obtém por um raciocínio que, ainda hoje, é um modelo de clareza
intelectual: o movimento do projétil, percebe ele, é na verdade
\textit{dois} movimentos sobrepostos e independentes. Na direção
horizontal, o projétil avança com velocidade constante --- nenhuma força
age sobre ele nessa direção, então nada o desacelera. Na direção
vertical, ele cai como qualquer corpo em queda livre --- acelerado pela
gravidade, percorrendo distâncias que crescem com o quadrado do tempo,
como as proporções dos números ímpares que ele mesmo havia descoberto
anos antes.

Dois movimentos simples. Combinados, produzem a elegância da parábola.

\begin{equation*}
    x = v_0\, t \qquad\qquad y = v_0\, t - \frac{1}{2}g t^2
\end{equation*}

O que a ilustração deste Primeiro Ato captura é precisamente esse
momento. Galileu está de costas para nós.
Não vemos seu rosto, não precisamos. O que ele vê é o que importa: as
posições sucessivas da bala, suspensas no ar como fantasmas de instantes
passados, traçando um arco que a geometria reconhece mesmo antes de a
física o explicar.

Sobre a mesa, o \textit{Discorsi e Dimostrazioni Matematiche} --- a obra
de 1638 na qual Galileu publica, já velho e cego, o resultado de décadas
de observação. Uma das principais referências usadas para construir o fio narrativo do primeiro ato desta obra. O compasso e a tinta ao lado: as ferramentas de quem mede
antes de concluir. A esfera armilar, que representa o cosmos antigo,
assiste em silêncio à inauguração de um novo modo de pensar.

E é esse novo modo que define o Primeiro Ato deste livro.

Antes de perguntar \textit{por que} os corpos se movem, é preciso
aprender a \textit{descrever} como eles se movem. Essa é a Cinemática
--- a linguagem do movimento puro, anterior às forças, anterior às
causas. Uma linguagem de posições, velocidades, trajetórias. Uma
linguagem que Galileu estava inventando naquele instante, olhando para
uma bala de canhão com a paciência de quem sabe que a natureza tem uma
resposta e está disposto a esperá-la.

A parábola que a bala traça no ar é descrita pelas mesmas equações que o leitor viu nas páginas deste Ato --- equações que Galileu estudou.
Ele não as escreveu exatamente assim --- essa notação veio depois. Mas o
pensamento é o mesmo: \textit{observar primeiro, medir e concluir
apenas quando a matemática permitir.}

Este Ato começa onde toda a Física começa --- com um homem olhando para
o mundo e tentando, com honestidade e rigor, dizer o que viu.







\section*{Segundo Ato --- Dinâmica}
% ─────────────────────────────────────────────────────────────────────

\begin{center}
    \textit{De Vi Invisibili} \\[0.2cm]
    {\small Sobre a força que não se vê --- e sobre aprender a explicar
    o que antes apenas se descrevia}
\end{center}

\bigskip

Há uma pergunta que Galileu não respondeu. Ele aprendeu a descrever o movimento com precisão extraordinária:
sabia dizer onde um corpo estaria, quão rápido iria, qual curva
desenharia no ar. Mas havia uma pergunta que ele deixou em aberto ---
talvez deliberadamente, talvez porque soubesse que ainda não era hora
de respondê-la:

\begin{center}
    \textit{Por que os corpos se movem como se movem?}
\end{center}

A resposta viria cinquenta anos depois, nas mãos de um homem sentado
à luz de uma vela, diante de uma folha em branco e de um universo
que ainda não sabia que estava prestes a ser explicado.

O homem é Isaac Newton, e a cena da ilustração deste Segundo Ato
captura o instante preciso dessa explicação. Sobre a mesa, o
\textit{Philosophiæ Naturalis Principia Mathematica} --- talvez o
livro mais importante já escrito em linguagem científica --- aberto
na página em que Newton demonstra algo que, até então, ninguém havia
sequer suspeitado: que a força que derruba uma maçã e a força que
mantém a Lua em órbita são \textit{a mesma força.}

O diagrama na página direita do livro mostra um canhão no topo de
uma montanha altíssima, acima da atmosfera. À medida que a
velocidade inicial do projétil aumenta --- $V_1, V_2, V_3, V_4$ ---
a trajetória deixa de ser uma parábola e começa a curvar em torno da
Terra. No limite, o projétil nunca chega ao chão: ele entra em órbita. A Lua,
percebe Newton, é exatamente isso --- um projétil que cai sem jamais
atingir o solo, porque a Terra curva sob seus pés na mesma taxa em
que ele desce.

Pela janela, a maçã pende do galho. O arco pontilhado a conecta à
Lua no céu noturno. Não é ornamento: é o argumento. A mesma
equação --- a mesma força, decaindo com o quadrado da distância ---
governa os dois movimentos. Um único princípio, dois fenômenos que
pareciam não ter nada em comum.

Mas o centro da ilustração não é o livro, nem a janela. É a folha
solta sobre a mesa, onde uma mão escreve, com tinta e compasso, o
diagrama de forças sobre uma esfera --- setas apontando para todos
os lados, convergindo para o centro --- e, abaixo, três símbolos que
mudariam a história:

\begin{equation*}
    \vec{F} = m\vec{a}.
\end{equation*}

Essa equação é o coração do Segundo Ato. Ela diz que a força não é
a causa do movimento --- como Aristóteles pensou e como se acreditou por dois mil
anos --- mas a causa da \textit{variação} do movimento. Um corpo em
repouso não precisa de força para permanecer em repouso. Um corpo em
movimento não precisa de força para continuar se movendo. O que
exige força é a mudança: acelerar, desacelerar, curvar.

É uma ruptura silenciosa. Não há drama na equação, nenhum anjo
descendo do céu, nenhuma revelação mística. Há apenas um homem, uma
vela, e a paciência de quem refaz os cálculos até que a natureza
concorde.

Esse é o espírito do Segundo Ato: aprender não apenas a descrever o
movimento, mas a \textit{explicá-lo.} A perguntar não só \textit{como,}
mas \textit{por quê.} A encontrar, por trás da trajetória que Galileu
mediu com tanto cuidado, a causa invisível que a produz.

A força não se vê. Mas seus efeitos, agora, têm nome e equação.




\section*{Terceiro Ato --- Aplicações}
% ─────────────────────────────────────────────────────────────────────
\begin{center}
    \textit{Harmonia Universi} \\[0.2cm]
    {\small Sobre a lei que governa planetas e navios --- e sobre o momento
    em que a teoria encontra o mundo}
\end{center}

\bigskip

Há uma pergunta que só pode ser feita depois. Não no início, quando ainda estamos aprendendo a linguagem. Não no
meio, quando ainda estamos construindo a teoria. Só no fim --- quando
a linguagem já foi aprendida e a teoria já está de pé --- é que se
pode olhar para o mundo com olhos novos e perguntar:

\begin{center}
    \textit{Isso aqui também obedece às mesmas leis?}
\end{center}

A resposta, invariavelmente, é sim. E é essa resposta que define o
Terceiro Ato.

A ilustração não tem personagem. Não há nenhum homem de costas
observando, nenhuma mão escrevendo à luz de vela. O que há é o
próprio mundo --- o cosmos acima e o oceano abaixo --- funcionando
em silêncio, sem precisar de testemunha. A Física, a essa altura,
já não depende de ninguém para existir. Ela simplesmente é.

No alto da imagem, o Sol irradia luz e, ao mesmo tempo, gravidade.
Os planetas se curvam ao seu redor em elipses que Kepler mediu e
Newton explicou. Saturno com seus anéis, Júpiter com sua massa
imensa, e ao fundo a faixa da Via Láctea lembrando que esse sistema
solar é apenas um ponto em uma escala ainda maior. No canto, um
cometa risca o céu --- o mesmo cometa que Halley, usando as equações
de Newton, ousou prever que voltaria. E voltou.


No centro, um navio. Ele não sabe que está provando um teorema.
Simplesmente flutua --- porque o peso do fluido deslocado pelo seu
casco é exatamente igual ao seu próprio peso, como Arquimedes
percebeu duzentos anos antes de Cristo numa banheira em Siracusa.
Abaixo da linha d'água, as linhas de pressão se irradiam para fora
e para baixo, em espelho perfeito das órbitas acima. A mesma
geometria, dois domínios completamente diferentes.


E nas rochas em primeiro plano, gravada como se sempre tivesse estado
ali, a equação que une os dois mundos:

\begin{equation*}
    F = G\,\frac{m_1 m_2}{r^2}
\end{equation*}

É a Lei da Gravitação Universal de Newton --- a mesma que curva os
planetas em órbita e a mesma que puxa o casco do navio em direção
ao fundo do mar, sendo contrariada, milímetro a milímetro, pela
pressão do fluido que o envolve. Não são duas forças diferentes.
É uma só, manifestando-se em escalas e contextos que Galileu e Newton
jamais imaginaram que um dia seriam explicados pela mesma sentença
matemática.

Esse é o privilégio do Terceiro Ato: chegar ao ponto em que a teoria
para de ser teoria e começa a ser realidade. Em que as equações que
você aprendeu a derivar no Primeiro Ato e a entender no Segundo Ato
revelam, de repente, que sempre estiveram ali --- nos planetas, nos
oceanos, nos navios, nos cometas que riscam o céu e voltam pontualmente,
como se soubessem que alguém os estava esperando.

A Física não descreve o mundo. Ela \textit{é} o mundo, escrito em
uma linguagem que levou séculos para ser decifrada --- e que você,
ao chegar até aqui, aprendeu a ler.




\fechamentodolivro